% Part: computability
% Chapter: computability-theory
% Section: complement-ce

\documentclass[../../../include/open-logic-section]{subfiles}

\begin{document}

\olfileid{cmp}{thy}{cmp}
\olsection{గణనీయంగా లెక్కించదగిన సమితులు పూరకం కింద సంవృతాలు కావు}

$A$ గణనీయంగా లెక్కించదగినదని అనుకుందాం. $A$ యొక్క పూరకం
$\Complement{A}=\Nat\setminus A$ కూడా ఎల్లప్పుడూ గణనీయంగా
లెక్కించదగినదేనా? సమాధానం ``కాదు'' అని కింది సిద్ధాంతం, దాని
ఉపసిద్ధాంతం చూపుతాయి.

\begin{thm}
\ollabel{thm:ce-comp}
$A$ సహజ సంఖ్యల ఏ సమితైనా అనుకుందాం. $A$, $\Complement{A}$ రెండూ
గణనీయంగా లెక్కించదగినవి అయితేనే $A$ గణనీయం.
\end{thm}

\begin{proof}
ముందు దిశ సులభం: $A$ గణనీయమైతే $\Complement{A}$ కూడా గణనీయం
($\Char{A}=1\tsub\Char{\Complement{A}}$); అందువల్ల రెండూ
గణనీయంగా లెక్కించదగినవే.

మరొక దిశలో, $A$,~$\Complement{A}$ రెండూ గణనీయంగా లెక్కించదగినవి
అనుకుందాం. $A$ అనేది~$\cfind{d}$ యొక్క నిర్వచన క్షేత్రం,
$\Complement{A}$ అనేది~$\cfind{e}$ యొక్క నిర్వచన క్షేత్రం అనుకుందాం.
$h$ను ఇలా నిర్వచించండి:
\[
h(x)=\umin{s}{(T(d,x,s)\lor T(e,x,s))}.
\]
మరో మాటలో, ఇన్‌పుట్~$x$పై $h$, $\cfind{d}$కు లేదా $\cfind{e}$కు
చెందిన ఆగే గణన కోసం వెతుకుతుంది. $x\in A$ అయితే మొదటి సందర్భంలో,
$x\in\Complement{A}$ అయితే రెండవ సందర్భంలో అది విజయవంతమవుతుంది.
కాబట్టి $h$ సర్వనిర్వచిత గణనీయ ప్రమేయం. ఇప్పుడు ప్రతి~$x$కు,
$x\in A$ అయితేనే $T(d,x,h(x))$; అంటే $\cfind{d}$ నిర్వచితమయ్యేదే
ఆ సందర్భం. $T(d,x,h(x))$ గణనీయ సంబంధం కనుక $A$ గణనీయం.
\sourcecorrection{OLTECOMTHY-008}{ఆంగ్ల మూలం $A$కు $\cfind{d}$ను, $\Complement{A}$కు $\cfind{e}$ను కేటాయించిన తరువాత కూడా, చివరి లక్షణీకరణలో పొరపాటుగా $T(e,x,h(x))$ అని రాసింది. $A$ను లక్షణీకరించడానికి అవసరమైన $T(d,x,h(x))$గా సరిచేశాం.}
\end{proof}

\begin{explain}
అనౌపచారిక గణనా పరంగా జరుగుతున్నది అర్థం చేసుకోవడం మరింత సులభం:
$A$ను నిర్ణయించడానికి, ఇన్‌పుట్~$x$పై $\cfind{d}$, $\cfind{e}$ల
ఆగే గణనల కోసం వెతకండి. వాటిలో ఒకటి తప్పక ఆగుతుంది; $\cfind{d}$
ఆగితే $x$ అనేది~$A$లో ఉంటుంది, లేకపోతే $x$ అనేది~$\Complement{A}$లో
ఉంటుంది.
\sourcecorrection{OLTECOMTHY-009}{ఈ వివరణలో ఆంగ్ల మూలం, నిరూపణలో నిర్వచించిన $d$,~$e$ సూచికలకు బదులు $e$,~$f$ అని రాసింది. వెంటనే ముందున్న కేటాయింపులకు అనుగుణంగా వాటిని $d$,~$e$గా సరిచేశాం.}
\end{explain}

\begin{cor}
\ollabel{cor:comp-k}
$\Complement{K_0}$ గణనీయంగా లెక్కించదగినది కాదు.
\end{cor}

\begin{proof}
$K_0$ గణనీయంగా లెక్కించదగినదే కాని గణనీయం కాదని మనకు తెలుసు.
$\Complement{K_0}$ కూడా గణనీయంగా లెక్కించదగినదైతే,
\olref{thm:ce-comp} ప్రకారం $K_0$ గణనీయం అవుతుంది; ఇది
\olref[ncp]{thm:K-0}కు వైరుధ్యం.
\end{proof}

\end{document}
